\chapter{Oral Candidiasis}
\index{Oral Candidiasis}

\medskip
\noindent\textit{\textbf{Under review.} This chapter is an AI-assisted educational draft; it is not a substitute for professional medical advice.}

\section*{Definition and Overview}
oral candidiasis is an overgrowth of Candida yeast in the mouth, often associated with antibiotics, inhaled corticosteroids, dentures, diabetes, or immune suppression.

\section*{Clinical Features}
white or red mouth patches, soreness, altered taste, cracked mouth corners, and discomfort when eating or swallowing.

\section*{When to Seek Medical Care}
Seek medical review for new, persistent, worsening, or function-limiting symptoms. Contact your local emergency services for severe breathing difficulty, collapse, sudden neurological change, severe chest or abdominal pain, or any rapidly worsening illness.

\section*{Diagnosis and Investigations}
oral examination usually establishes the diagnosis; recurrent or severe disease may require review for medicines, diabetes, immune problems, or an alternative diagnosis.

\section*{Management}
topical or oral antifungal treatment is selected by a clinician; rinse the mouth after inhaled steroids, clean dentures, and address predisposing factors.

\section*{Complications}
Complications depend on the underlying cause and may include progression, effects on other organs, treatment adverse effects, or reduced quality of life.

\section*{Prognosis and Self-Care}
Outcome depends on the cause, severity, complications, and response to treatment. Follow the care plan, keep appointments, avoid known triggers, and seek help if symptoms worsen.
\section*{Safety-netting}
Seek urgent care for sudden severe symptoms, chest pain, breathing difficulty, collapse, new neurological deficit, or rapidly worsening function.
